\section{A positive route: intervention monotonicity}\label{sec:via-futura}
Theorem~\ref{thm:dicotomia} closes one door. If equality of degree is required on every pair of $\ker\Phi_Q$, and the defended feature $\kappa$ co-occurs connectedly with the levels of $\Phi_Q$, the degree is constant. This section asks which reformulation of $S_u$ escapes that collapse without circularity. Two natural candidates do not escape (Section~\ref{sec:collapse}). One does (Section~\ref{sec:im}), at a price that we state exactly; we have not shown that it is the only reformulation that does. The surviving route can be instantiated in the band of Act~I through a monotonicity theorem for $J$, proved on paper and checked in Lean with Mathlib (Section~\ref{sec:scaling}). We close with a methodological guide and with the statements that should be formalized next. The Lean statements about degrees cited in this section are in \texttt{ViaFutura.lean} (the graph lemmas cited in Corollary~\ref{cor:band-scaling} are in \texttt{PistasUnificacion.lean}); they take degrees with integer values, and their arguments use only the structure of a totally ordered abelian group, so the statements carry over verbatim to real values, with the same arguments; that transcription is not formalized, and the Lean proofs, which use integer tactics such as \texttt{omega} and \texttt{decide}, would have to be rewritten. The statements of Proposition~\ref{prop:scaling} are checked with Mathlib, over real matrices, in \texttt{Scaling.lean}, \texttt{Interval.lean} and \texttt{Bridge.lean}.

\subsection{Two reformulations that collapse}\label{sec:collapse}
\paragraph{Exact covariance under perturbations $\Delta\Phi_Q$.} Let $\mathcal T$ be a class of perturbation steps, and suppose that along every step $u$ changes exactly as a function of $\Phi_Q$: $u(t)-u(s)=H(\Phi_Q(t))-H(\Phi_Q(s))$ for $(s,t)\in\mathcal T$. Then $u-H\circ\Phi_Q$ is constant along chains of steps, so the two ends of any directed chain with equal $\Phi_Q$ have equal $u$ (\texttt{exact\_covariance\_gives\_Su}) \stL. Exact covariance is symmetric, so the same holds along the symmetric closure of $\mathcal T$, and $S_u$ holds on every connected component of the undirected step graph (\texttt{exact\_covariance\_components}) \stL. Exact covariance is therefore $S_u$ on components, not an escape from it.

\paragraph{The order clause B on all pairs.} Suppose $u=f\circ\kappa$ and B holds on all pairs: $\Phi_Q(s)<\Phi_Q(t)$ implies $f(\kappa(s))\le f(\kappa(t))$. If a class $k$ occurs strictly below a class $k'$ and $k'$ occurs strictly below $k$, then $f(k)=f(k')$ (\texttt{order\_all\_pairs\_general}) \stL. With a transversal $\kappa$ whose image is a product of at least two distinct levels of $\Phi_Q$ and the realized classes, every pair of realized classes is in this situation, and $f$ is constant on them; \texttt{order\_all\_pairs\_collapse} \stL\ is the instance with two levels and two classes.

\paragraph{A cohomological obstruction $H^1$ to gluing bands of SEG is ill-posed.} Consider the presheaf of local factorizations, $\mathcal F(U)=\{h:U\cap\operatorname{im}\Phi_Q\to\mathbb R \mid u=h\circ\Phi_Q \text{ on } \Phi_Q^{-1}(U)\}$. It is a sheaf, and its sections over any $U$ are either empty or unique. SEG, and with it $S_u$, is a condition on stalks (local existence of a factorization), not a gluing condition, so $H^1$ detects nothing. With coefficients in calibrations, such as monotone rescalings of $u$, the parameter space of stable linear-Gaussian systems is star-shaped ($tA$ is stable for $t\in[0,1]$ whenever $A$ is), hence contractible, and the constant-coefficient $H^1$ vanishes as well. \stA

\subsection{The route that survives}\label{sec:im}
Let $\mathcal T$ be a declared class of interventions on one and the same bearer, for instance deformations of its couplings. \emph{Intervention monotonicity} is the requirement
\begin{equation}\label{eq:im}
\mathrm{IM}_{\mathcal T}:\qquad (s,t)\in\mathcal T,\ \ \Phi_Q(s)<\Phi_Q(t)\ \Longrightarrow\ u(s)\le u(t).
\end{equation}
For $u=f\circ\kappa$, the \emph{intervention digraph} has the realized coexperience classes, $\operatorname{im}\kappa$, as vertices, and an edge $\kappa(s)\to\kappa(t)$ for every $(s,t)\in\mathcal T$ with $\Phi_Q(s)<\Phi_Q(t)$.
\begin{theorem}[Exact frontier; \texttt{im\_iff\_edges}, \texttt{constant\_on\_scc}, \texttt{nonconstant\_if\_not\_scc}]\label{thm:frontier}
Let $u=f\circ\kappa$.
\begin{enumerate}
\item[(a)] $\mathrm{IM}_{\mathcal T}$ holds if and only if $f$ is monotone along the edges of the intervention digraph.
\item[(b)] Every such monotone $f$ is constant on each strongly connected component of the digraph.
\item[(c)] If $k=\kappa(s)$ and $k'=\kappa(s')$ are realized classes and $k$ is not reachable from $k'$, there is an $f$, monotone along the edges, with $f(k)\ne f(k')$, and hence $u(s)\ne u(s')$.
\end{enumerate}
All three parts are compiled in Lean. \stL
\end{theorem}
The route therefore escapes the collapse exactly when the intervention digraph on the realized classes is not strongly connected: when the interventions that raise $\Phi_Q$ cannot carry every realized coexperience class to every other one in both directions. A finite model escapes with a connected product image (\texttt{intervention\_monotonicity\_escapes}) \stL. The frontier is sharp in the other direction too: if equality is also required on $\Phi_Q$-preserving steps, and such steps chain together the realizations of a fibre, $S_u$ returns on that fibre (\texttt{equality\_on\_fibre\_steps\_collapses}) \stL.

\emph{Status.} $\mathrm{IM}_{\mathcal T}$ is \emph{not} $S_u$, and it is not original Q. It compares only states of one bearer that are related by declared interventions, and it compares them ordinally. It is a declared weakening of the order clause B to a class of interventions. In exchange, it needs neither a selection of representatives nor SEG, and it asks only for ordinal comparisons within a subject, not for absolute comparisons between different architectures.

\subsection{An instance in the band: monotonicity of \texorpdfstring{$J$}{J} under scaling}\label{sec:scaling}
\begin{proposition}[Monotonicity under scaling]\label{prop:scaling}
Let $n\ge2$ and $E\succ0$.
\begin{enumerate}
\item[(a)] \emph{Uniform scaling.} Let $\operatorname{spr}(A)<1$ and $t\in[0,1/\operatorname{spr}(A))$ (all $t\ge0$ if $\operatorname{spr}(A)=0$). Then $t\mapsto J(tA,E)$ is nondecreasing, with $J(0,E)=0$, and it is strictly increasing if the undirected support of $A$ is connected.
\item[(b)] \emph{Fixed scalar self-coupling.} Let $A(c)=aI+cW$ with $|a|<1$ and $W$ of zero diagonal, and let $c$ range over $\mathcal C=\{c\ge0:\operatorname{spr}(A(c))<1\}$, an interval containing $0$ (each eigenvalue $a+c\mu$ of $A(c)$ has $|a+c\mu|$ convex in $c$). Then $c\mapsto J(A(c),E)$ is nondecreasing on $\mathcal C$, and it is strictly increasing if the undirected support of $W$ is connected.
\end{enumerate}
In both cases, for $t>0$ and $c>0$ the support is that of $A$ (respectively $W$), and in a fixed representative the minimum cut grows as $t^2$ (respectively $c^2$). \stA\ \stM
\end{proposition}
\begin{proof}
(a) (1) $P_t=\sum_{j\ge0}t^{2j}A^jE(A^j)^\top$ solves $P_t=(tA)P_t(tA)^\top+E$, and $P_{t_2}-P_{t_1}=\sum_j(t_2^{2j}-t_1^{2j})A^jE(A^j)^\top\succeq0$ for $t_1\le t_2$: $P_t$ is nondecreasing in the Loewner order.
(2) Fix a bipartition $\pi$ and $t>0$. The map $B'\mapsto tB'$ is a bijection of $\mathfrak B_\pi$, and $E+(tA-tB')P_t(tA-tB')^\top=E+t^2(A-B')P_t(A-B')^\top=:R_t(B')$. For fixed $B'$ and $0<t_1<t_2$, with $C=A-B'$,
\begin{equation}\label{eq:resid-incr}
R_{t_2}(B')-R_{t_1}(B')=t_2^2\,C(P_{t_2}-P_{t_1})C^\top+(t_2^2-t_1^2)\,CP_{t_1}C^\top\succeq0,
\end{equation}
by (1) and because $X\mapsto CXC^\top$ preserves the Loewner order. Since $\log\det$ is monotone in the Loewner order, $f_t(B')=\tfrac12\log\det R_t(B')-\tfrac12\log\det E$ is nondecreasing in $t$, and so are its infimum over $B'$ and the minimum over $\pi$. At $t=0$, $B=0$ gives $J=0$, and $J\ge0$ always.
(3) If the support is connected, $A\notin\mathfrak B_\pi$ for every $\pi$, so $A-B'\ne0$ for every $B'\in\mathfrak B_\pi$. Let $0<t_1<t_2$, and let $B'_*$ minimize $f_{t_2}$ on $\mathfrak B_\pi$. It exists because $f_{t_2}$ is continuous and coercive: $f_t(B')\ge\tfrac12\log\bigl(1+t^2\lambda_{\min}(P_t)\lambda_{\max}(E)^{-1}\lVert A-B'\rVert_F^2\bigr)$, using $\det(I+X)\ge1+\operatorname{tr}X$ for $X\succeq0$. Put $M=t_1^2(A-B'_*)P_{t_1}(A-B'_*)^\top$, which is positive semidefinite and nonzero because $P_{t_1}\succ0$. By (1), $R_{t_2}(B'_*)\succeq E+(t_2/t_1)^2M$, and $\lambda\mapsto\log\det(E+\lambda M)$ is strictly increasing on $\lambda\ge0$, with derivative $\operatorname{tr}((E+\lambda M)^{-1}M)>0$. Hence, writing $J_\pi(tA)=\min_{B'\in\mathfrak B_\pi}f_t(B')$ for the inner minimum, $J_\pi(t_2A)=f_{t_2}(B'_*)>f_{t_1}(B'_*)\ge J_\pi(t_1A)$ for every $\pi$; taking $\pi$ minimizing at $t_2$ gives $J(t_2A)>J(t_1A)$. For $t_1=0$, $J(0)=0<J(t_2A)$ by Lemma~\ref{lem:J}.
(4) For $t>0$ the nonzero coefficients of $tA$ are those of $A$ multiplied by $t$; the support is unchanged and the crossing coefficients scale by $t$.

(b) Let $P(c)$ be the stationary covariance of $A(c)$ and $Q(c)=c^2P(c)$ (a matrix, unrelated to the theory Q), so that $Q=AQA^\top+c^2E$. Since $aI\in\mathfrak B_\pi$, for $c>0$ the map $B'\mapsto aI+cB'$ is a bijection of $\mathfrak B_\pi$, and $A(c)-(aI+cB')=c(W-B')$; the residual is $E+(W-B')Q(c)(W-B')^\top$. Differentiating $Q=AQA^\top+c^2E$ gives $Q'=AQ'A^\top+S$ with $S=WQA^\top+AQW^\top+2cE$, and the Lyapunov equation for $P$ turns this into
\begin{equation}\label{eq:S-identity}
S=c\bigl[(1-a^2)P+c^2WPW^\top+E\bigr]\succ0\qquad(c>0).
\end{equation}
Hence $Q'=\sum_{j\ge0}A^jS(A^j)^\top\succ0$, and $Q(c)$ is strictly increasing in the Loewner order. The rest of the argument is that of (a), with $Q(c)$ in place of $t^2P_t$ and $W-B'$ in place of $A-B'$: for fixed $B'$ the residual is nondecreasing in $c$, and for $0<c_1<c_2$ the residual at $c_2$ exceeds the residual at $c_1$ by $(W-B')\bigl(Q(c_2)-Q(c_1)\bigr)(W-B')^\top$, which is positive semidefinite and, when the support of $W$ is connected (so that $W-B'\ne0$ for every $B'\in\mathfrak B_\pi$), nonzero. The minimizer exists by coercivity, since $Q(c)\succ0$ for $c>0$. At $c=0$ the support is empty and $J=0$.
\end{proof}
{\sloppy \emph{Formal check} \stM. Both parts of the proposition are checked in Lean with Mathlib, over real matrices, as \texttt{QAuditScaling.prop\_scaling\_a} and \texttt{QAuditScaling.prop\_scaling\_b}; the other theorems named in this paragraph are in the same namespace, \texttt{QAuditScaling}, unless another one is given. The last sentence of the proposition is \texttt{supportGraph\_smul} and \texttt{supportGraph\_Acoup} for the support, and \texttt{wmin\_smul} and \texttt{wmin\_Acoup} for the minimum cut. The stationary covariance of $M=tA$ or $M=A(c)$ is the series $\sum_{j\ge0}M^jE(M^j)^\top$, which solves $P=MPM^\top+E$ (\texttt{Pt\_lyap}, \texttt{Pser\_lyap}) and whose convergence is proved from $\operatorname{spr}(M)<1$ by Gelfand's formula (\texttt{serSummable\_of\_spr\_lt\_one}). Here $\operatorname{spr}$ is Mathlib's spectral radius of $M$ regarded as a complex matrix; that it is the largest modulus of an eigenvalue is standard and is not formalized. The range of $t$ in (a) is written as $t\operatorname{spr}(A)<1$, using $\operatorname{spr}(tA)=t\operatorname{spr}(A)$ (\texttt{spr\_smul}). The inner minimum is an infimum, and it is attained (\texttt{exists\_minimizer}). The checked statement is slightly more general: (a) does not assume $\operatorname{spr}(A)<1$, and (b) does not assume that $W$ has zero diagonal. Neither $Q'=AQ'A^\top+S$ nor the identity~\eqref{eq:S-identity} is stated in Lean: the Lean proof of (b) replaces the derivative by the inequality $Q(c_2)-Q(c_1)\succeq(c_2-c_1)^2E$ for $c_1\le c_2$ in $\mathcal C$ (\texttt{Qb\_mono}). That $\mathcal C$ is an interval containing $0$ is checked with the argument given in (b), by \texttt{C\_ordConnected} and \texttt{C\_interval\_containing\_zero} in the namespace \texttt{QAuditScaling.AuditInterval}. \texttt{Scaling.lean} defines $\mathfrak B_\pi$, $J$, $w_{\min}$ and the support graph independently of the Lean statements of Lemma~\ref{lem:J}, and \texttt{QAuditBridge.J\_eq}, \texttt{QAuditBridge.wmin\_eq} and \texttt{QAuditBridge.supportGraph\_eq} prove that the two sets of definitions agree, those of $J$ for $E\succ0$ and $P\succeq0$.\par}

The identity~\eqref{eq:S-identity} is an exact polynomial identity modulo the Lyapunov equation, checked with generic symbolic $3\times3$ matrices (for general $n$ it is the computation in the proof \stA), and $Q'=AQ'A^\top+S$ holds exactly on a rational instance with symbolic $c$ (\texttt{escalado\_diagonal\_escalar.py}) \stS.

\begin{remark}[The diagonal must scale, or be scalar]\label{rem:fixed-diagonal}
Proposition~\ref{prop:scaling} does not extend to strengthening the couplings while a \emph{non-scalar} diagonal of self-couplings is held fixed; then $S$ can be indefinite. With $n=2$, $A(c)=D+cW$, $D=\mathrm{diag}(-0.23,\,0.31)$, $W=\bigl(\begin{smallmatrix}0&-0.19\\-0.83&0\end{smallmatrix}\bigr)$ and $E=\bigl(\begin{smallmatrix}4.61&2.59\\2.59&1.46\end{smallmatrix}\bigr)\succ0$, one has $J(A(0.33))\approx0.104083>J(A(0.38))\approx0.099702$ (\texttt{contraejemplo\_diagonal\_fija.py}: exact elimination of one variable and real roots to 50 digits) \stN. The four-unit band systems of Table~\ref{tab:networks}, with scalar self-coupling $a=0.2$ and all couplings equal to $c$, are covered by part~(b).
\end{remark}
\begin{corollary}[The band under scaling]\label{cor:band-scaling}
For every $n\ge4$, a stable dynamics whose undirected off-diagonal support is exactly the path $P_n$ (in~(b), $A=aI+cW$ with the undirected support of $W$ exactly $P_n$) stays in the band along the deformations of Proposition~\ref{prop:scaling} with $t>0$ (respectively $c>0$): its support is unchanged, so it remains connected without a universal vertex (\texttt{path\_reach\_from\_zero}, \texttt{path\_no\_universal}) \stL. Along these deformations $J$ increases strictly and the minimum cut grows quadratically \stA\ \stM. The same holds when the support is exactly the cycle $C_n$, $n\ge4$, whose maximum degree $2$ is smaller than $n-1$.
\end{corollary}
The corollary combines two checks, the graph part without Mathlib, on the Boolean adjacency of $P_n$, and the part on $J$ with Mathlib, through Proposition~\ref{prop:scaling}; it is not a single Lean statement.

Numerically, $J$ increases strictly on grids of $t$ (uniform scaling) for $P_4$, $P_5$, $P_6$, $C_4$, $C_5$ and five random connected systems with random $E$, with $w_{\min}(tA)=t^2w_{\min}(A)$ to machine precision (\texttt{monotonia\_escalado.py}); and on grids of $c$ (scalar self-coupling $a=0.2$) for $P_4$, $C_4$, $P_5$, $C_5$ and four random connected systems (\texttt{escalado\_diagonal\_escalar.py}, \texttt{banda\_monotonia.py}) \stN.

\emph{What $\mathrm{IM}_{\mathcal T}$ would then predict.} Take $\mathcal T$ to be one of these deformations of a band system, increasing in $t$ or in $c$. If $J$ were taken as the descriptor in $\mathrm{IM}_{\mathcal T}$, which this paper does not justify because $J$ is not $\Phi_Q$ (Section~\ref{sec:gauss}), $\mathrm{IM}_{\mathcal T}$ would predict that $u$ does not decrease as the couplings of the system are strengthened in that way. With the original descriptor the prediction also requires that $\Phi_Q$ increase along the deformation, which has not been shown. Every edge of the corresponding intervention digraph goes from a smaller to a larger parameter. The digraph on the realized classes is strongly connected only if coexperience classes recur in both directions along the deformation, and whether they do is an empirical, intra-subject question. Testing the prediction needs an ordinal intra-subject measure of $u$ along such interventions, that is, a calibrated measurement channel for $u$, which is not available: the stipulation A14--U of Definition~\ref{def:orderB} fixes a rule assigning observables to the degree, not a calibration.

\subsection{A methodological guide}\label{sec:guide}
The results of the three acts suggest a way of comparing theories of consciousness that neither presupposes $S_u$ nor collapses it.
\begin{enumerate}
\item \emph{Do not equate degrees across different architectures.} Static equality across systems with equal descriptor is exactly the obligation that the principles examined in Act~I cannot discharge without circularity, collapse, or a change of the meaning of the degree or of the domain.
\item \emph{Locate the disagreement first.} The structural presence bridges of IIT type and GNWT type commute on the central and disconnected levels and disagree exactly on the band (Section~\ref{sec:l1}). Systems in the band, such as $P_n$ and $C_n$ for $n\ge4$, are where the two kinds of proxy make different predictions.
\item \emph{Intervene within one system.} Use declared interventions, such as the deformations of Proposition~\ref{prop:scaling}, and ask only for ordinal comparisons of $u$ along them.
\item \emph{Check the digraph.} By Theorem~\ref{thm:frontier}, the comparison is informative exactly outside the strongly connected components of the intervention digraph on the realized classes; within a component, $u$ is forced to be constant.
\end{enumerate}

\subsection{What should be formalized next}\label{sec:next}
\begin{enumerate}
\item[(i)] In Lean~4 with Mathlib, the analytic statements that are not yet formalized: the hub bound of Theorem~\ref{thm:laxa}(d), the decomposition $I_{\rm pred}=I_{\rm cut}+J$ and the additivity of $I_{\rm pred}$ under $\otimes$ (Proposition~\ref{prop:monoidal}), which are paper proofs, and the statements about the restricted Gaussian family of Section~\ref{sec:gauss}, which are exact computations in SymPy or paper proofs (only the value $J=\tfrac12\log\tfrac54$ on the member $\theta=(\tfrac12,1,1)$ is checked with Mathlib, \texttt{QAudit.Audit.J\_ex\_eq}). Lemma~\ref{lem:J}, the min-cut sandwich (Theorem~\ref{thm:cheeger}) and Proposition~\ref{prop:scaling} are already checked with Mathlib \stM; what links them to the graph facts checked without Mathlib, such as the connectivity of $P_n$ in Corollary~\ref{cor:band-scaling} and the disconnected support of a tensor product, is still on paper, and so are the conventions of Remark~\ref{rem:small} for fewer than two units.
\item[(ii)] For a physically specified class $\mathcal T$ of deformations on $\mathbf{Open}_G$, the structure of strongly connected components of its intervention digraph on coexperience classes. It is the exact condition under which the route does not degenerate.
\item[(iii)] Not a theorem but its empirical prerequisite: a calibrated, ordinal, intra-subject measurement link for $u$, which the stipulation A14--U does not supply (Definition~\ref{def:orderB}(c)).
\end{enumerate}
Combining the zero-product monoid and the additive monoid of Section~\ref{sec:l2} into one structure with two operations only reorganizes the bookkeeping, and using $P_n$ as a test bench requires data on $u$ that do not exist; neither of these alternatives addresses $S_u$.
